\chapter{The database}\label{sec:db}
\section{Contents and construction}
The package \lcodb{} reads the LCO science archive (\texttt{archive-api.lco.global}) and the observation portal
(\texttt{observe.lco.global}) with the PI's API token and writes two SQLite databases:
\begin{itemize}\setlength{\itemsep}{1pt}
\item \texttt{photometry.sqlite}: Sinistro (camera prefixes \texttt{fa}, and \texttt{fl} for the early
  Sinistros), Spectral (\texttt{fs}) and MuSCAT (\texttt{ep}) images. Only reduced frames are kept: the BANZAI
  product (reduction level 91) and, where none exists, the earlier pipeline product (level 90; flag
  \texttt{is\_best}).
\item \texttt{spectra.sqlite}: FLOYDS raw frames (science, arcs, lamp flats) and BANZAI-FLOYDS products,
  plus the flux standards from any proposal (\S\ref{sec:standards}).
\end{itemize}
NRES, guide and acquisition cameras, catalogue files and experimental frames are excluded. Each frame keeps its
archive identifiers (so the file can be fetched later), target, coordinates, filter, exposure time, observing
block and request, the start time as ISO date, readable UTC (\texttt{date\_utc}) and MJD, and a set of header
keywords read from the archive's header service, which returns the full FITS header without the pixels. The
full header is also stored (compressed), so further keywords can be added as columns without new downloads.
An airmass is computed for every frame from the site, the footprint centre and the mid-exposure time; it agrees
with the header \texttt{AIRMASS} to 0.01.

\section{Header keywords across the decade}
A sample of headers from every year, camera and reduction level (2014--2026) showed that the keywords needed
for basic image quality are not homogeneous:
\begin{itemize}\setlength{\itemsep}{1pt}
\item \texttt{MJD-OBS} and \texttt{DATE-OBS} are missing from the headers served for frames before 2023; the
  database takes the time from the archive record instead, which is complete.
\item \texttt{TELID} only gives the class (\texttt{1m0a}); the individual telescope comes from the file name
  (\texttt{1m010}).
\item \texttt{SKYMAG} is the same in raw and reduced headers, i.e.\ it is a telemetry/scheduler value, and
  takes the placeholder 22.0 in many frames (\fracSkyPlaceholder\,\% of the 1-m images), which the quality view
  treats as missing.
\item \texttt{L1FWHM} uses $-99$ for failures in the old pipeline; \texttt{NAXIS1/2} is reliable in the
  reduced frames (4096 for BANZAI Sinistro, $4018\times3974$ for the old Sinistro products, 2009--2045 for
  binned Spectral, 2048 for MuSCAT).
\end{itemize}

\section{Daily update}
The first build crawled each proposal from its first semester in 30-day windows. The daily run re-reads the last
seven days of every active proposal, so late frames and new BANZAI versions are picked up (a changed checksum
triggers a new header download), fetches the missing headers, and rebuilds the calibration and standard
associations. New proposals in the PI profile are added automatically.
